\section{Conclusion}
% PAPER-APP[claude:PR-STYLE-0004,codex:PR-STYLE-0002] pass=style state=open
% Issue (claude): Vague evaluative phrases that state a quality without saying how it is measured.
% Evidence (claude): here "semantically simple and easy to implement" and "great expressive power"; Section 5.1
%   "The evaluator can be simplified greatly".
% Proposal (claude): Replace with concrete properties already shown in the paper, e.g. "three special forms and the four
%   rule groups of the small-step semantics", "modular programs with named functions, mutual recursion and imports, none
%   built into the base language", "the evaluator distinguishes only special forms and functions".
% Reason (claude): Concrete statements are more persuasive to PEPM reviewers than adjectives; whether the claims
%   are supported is handed to the technical pass.
% Issue (codex): The conclusion relies on broad evaluative wording instead of naming the demonstrated capabilities.
% Evidence (codex): The paragraph describes ILD as "semantically simple", "easy to implement", and having "great expressive power".
% Proposal (codex): Replace these descriptions with a concise account of the three core special forms and the library constructions for binding, recursion and module loading; retain only claims established in the paper.
% Reason (codex): Concrete capabilities give the reader a precise statement of the outcome.
% END-PAPER-APP[claude:PR-STYLE-0004,codex:PR-STYLE-0002]
% PAPER-APP[claude:PR-STRUCTURE-0007] pass=structure state=open
% Issue: The conclusion does not mirror the abstract or answer the problem posed in the introduction.
% Evidence: The introduction motivates modular programs without mutation, and the abstract claims a module system
%   (named functions, mutual recursion, imports) and a hygiene benefit. The conclusion mentions neither the module
%   loader nor hygiene, and introduces "as powerful as traditional macro systems", which the body argues only informally.
% Proposal: Restructure as: problem (mutation-free modular programs) -> approach (ILD, callsite-triggered macros) ->
%   what was shown (library up to letrec, module loader without built-in features, hygiene through embedded values)
%   -> limitations and next step (performance, partial evaluation).
% Reason: Readers who skim abstract and conclusion should see the same contributions.
% END-PAPER-APP[claude:PR-STRUCTURE-0007]
We set out to support modular
programs without mutation of a global environment, to which end we present the
programming language ILD. Its semantics (summarised in \cref{stepped-semantics-fig})
is minimal: the only possible \emph{combination} steps are function application and
one of the three special forms. Together with a set of host functions
and data structures, provided via FFI, and a metaprogramming library written in ILD itself,
it lets the user write modular programs
with named functions, mutual recursion and source imports, none of which
is built into the base language. Callsite-triggered macro expansion,
where macros are ordinary first-class functions called at runtime,
allows for a simple evaluator with no separate macro expansion phase.
The macros in this paper make use of this runtime semantics to
emit values instead of keywords, which avoids name capture and leakage without
hygiene machinery.
In the future we aim to improve performance
by employing partial evaluation as an optimisation step.
